\documentclass[11pt,a4paper]{article}
\usepackage[utf8]{inputenc}
\usepackage[T1]{fontenc}
\usepackage[ngerman]{babel}
\usepackage[margin=0.75in]{geometry}
\usepackage{enumitem}
\usepackage{titlesec}
\usepackage{xcolor}
\usepackage{hyperref}
\definecolor{primarycolor}{RGB}{33,37,41}
\definecolor{secondarycolor}{RGB}{0,102,204}
\definecolor{accentcolor}{RGB}{108,117,125}
\hypersetup{colorlinks=true,linkcolor=secondarycolor,urlcolor=secondarycolor,citecolor=secondarycolor}
\titleformat{\section}{\large\bfseries\color{primarycolor}}{}{0em}{}
\titlespacing*{\section}{0pt}{1em}{0.4em}
\newcommand{\cvsection}[1]{\section*{#1}}
\newcommand{\cvitem}[2]{\noindent\textbf{#1}\\#2\par}
\setlist[itemize]{leftmargin=*,itemsep=0.15em,topsep=0.25em,parsep=0pt}
\setlength{\parindent}{0pt}
\setlength{\parskip}{0.25em}
\pagestyle{empty}
\begin{document}
{\LARGE\bfseries\color{primarycolor} KATARZYNA WIŚNIEWSKA}\\[0.3em]
{\large\color{accentcolor} Pflegehelferin}\\[0.3em]
Habsburgerstraße 41, 79104 Freiburg im Breisgau\\
k.wisniewska@example.com \textbar{} +49 761 555 0412

\cvsection{PROFIL}
Pflegehelferin mit Erfahrung in der stationären Altenpflege seit April 2020 und vorheriger Tätigkeit als Krankenschwester in Polen. Sicher in der Grundpflege, Behandlungspflege nach Leistungsgruppe 1 und digitalen Pflegedokumentation mit Vivendi PD. Ruhig, zuverlässig und geduldig im Umgang mit Menschen mit Demenz. Deutschkenntnisse auf Niveau B2.

\cvsection{BERUFSERFAHRUNG}
\cvitem{Pflegehelferin \textbar{} Seniorenzentrum Am Lindenhof, Freiburg}{04/2020 – heute}
\begin{itemize}
\item Grundpflege für 10 bis 12 Bewohnerinnen und Bewohner pro Schicht: Körperpflege, Mobilisation, Lagerung und Hilfe beim Essen.
\item Behandlungspflege nach Leistungsgruppe 1 (LG1): Blutdruck- und Blutzuckermessung, Kompressionsstrümpfe sowie Medikamente nach Plan reichen.
\item Digitale Pflegedokumentation mit Vivendi PD.
\item Erfahrung in Früh-, Spät- und Wochenenddiensten sowie in der Betreuung von Menschen mit Demenz.
\item Einarbeitung von bisher sieben neuen Pflegehelferinnen und Praktikanten während ihrer ersten zwei Wochen nach dem Einarbeitungsplan der Einrichtung.
\end{itemize}

\cvitem{Aushilfe Verkauf \textbar{} Bäckerei Schwarzwaldkorn, Freiburg}{09/2019 – 03/2020}
\begin{itemize}
\item Verkauf, Kassenbedienung und Reinigung.
\end{itemize}

\cvitem{Krankenschwester (Pielęgniarka) \textbar{} Wojewódzki Szpital Specjalistyczny, Opole, Polen}{10/2014 – 06/2019}
\begin{itemize}
\item Grund- und Behandlungspflege auf einer Station für Innere Medizin mit 32 Betten; Infusionen und Vorbereitung von Visiten.
\item Verbandswechsel bei Dekubitus und diabetischem Fuß unter Aufsicht der Stationsärztin, 2016–2019.
\end{itemize}

\newpage
\cvsection{AUSGEWÄHLTE PRAXISPROJEKTE}
\textbf{Sturzprophylaxe \textbar{} Seniorenzentrum Am Lindenhof \textbar{} 2023}
\begin{itemize}
\item Gemeinsam mit der Wohnbereichsleitung Sturzprotokolle aus sechs Monaten ausgewertet; Nachtlichter und feste Toilettengänge für neun Bewohner eingeführt.
\item Im folgenden Halbjahr sank die Zahl der Stürze im Wohnbereich von 14 auf 8.
\end{itemize}

\textbf{Nachmittagsgruppe für Menschen mit Demenz \textbar{} Seniorenzentrum Am Lindenhof \textbar{} seit 2022}
\begin{itemize}
\item Leitung einer wöchentlichen Singgruppe mit polnischen und deutschen Volksliedern für sechs bis acht Bewohner; positive Rückmeldungen von Angehörigen im Heimbeirat.
\end{itemize}

\cvsection{AUSBILDUNG UND ANERKENNUNG}
\textbf{Bachelor Pielęgniarstwo (Krankenpflege)}\\
Medizinische Universität Opole, 2014\\
Anerkennungsverfahren als Pflegefachfrau beim Regierungspräsidium Stuttgart läuft; Defizitbescheid liegt vor. Kenntnisprüfung für März 2027 vorgesehen. Derzeit in Deutschland als Pflegehelferin tätig, noch keine Anerkennung als Pflegefachfrau.

\cvsection{WEITERBILDUNG}
\begin{itemize}
\item Erste-Hilfe-Kurs, 2024
\item Expertenstandard Sturzprophylaxe, interne Schulung, 2023
\item Behandlungspflege LG1, Basiskurs, 2021
\end{itemize}

\cvsection{KENNTNISSE UND SPRACHEN}
Pflegedokumentation: gute Kenntnisse in Vivendi PD.\\
Deutsch: B2, telc Deutsch B1-B2 Pflege, Zertifikat 2021.\\
Polnisch: Muttersprache.\\
Englisch: Grundkenntnisse.

\cvsection{ARBEITSZEIT UND FÜHRERSCHEIN}
Gewünschter Umfang: Teilzeit mit 75 \%; Dienste an jedem zweiten Wochenende möglich.\\
Führerschein Klasse B noch nicht vorhanden; Fahrschulausbildung läuft, Theorieprüfung bestanden, praktische Prüfung für Januar 2027 vorgesehen.
\end{document}